% Original analytic teaching diagram, 2026-11-09.
\definecolor{VB}{HTML}{7188A3}\definecolor{VBF}{HTML}{E5EAF0}
\definecolor{VG}{HTML}{829C8B}\definecolor{VGF}{HTML}{E8EEE9}
\definecolor{VY}{HTML}{B4A16F}\definecolor{VYF}{HTML}{F1ECDD}
\definecolor{VR}{HTML}{A7787D}\definecolor{VRF}{HTML}{F0E5E6}
\definecolor{VP}{HTML}{9588A5}\definecolor{VPF}{HTML}{ECE8F0}
\definecolor{VI}{HTML}{3B4252}\definecolor{VM}{HTML}{6B7280}\definecolor{VL}{HTML}{4C4D4F}
\begin{tikzpicture}[x=1mm,y=1mm,>=Stealth,font=\figurefont\fontsize{8}{10}\selectfont,
lab/.style={text=FigureMuted,font=\figurefont\fontsize{7.5}{9}\selectfont},
box/.style={draw=VB,fill=VBF,line width=.55pt,rounded corners=1mm,minimum width=22mm,minimum height=12mm,align=center,text=VI},
fin/.style={draw=VL,line width=.45pt,rounded corners=1mm,minimum width=22mm,minimum height=9mm,align=center,text=VI},
ar/.style={->,draw=VL,line width=.6pt}]
\node[fin] (s0) at (0,30) {极小化序列\\$J(u^{(k)})\to J^\star$};
\node[box] (a) at (0,12) {序列有界\\强制性或\\有界下水平集};
\node[box] (b) at (26,12) {弱收敛子列\\反射性};
\node[box] (c) at (52,12) {弱极限可行\\$\bm{C}$ 弱序列闭};
\node[box] (d) at (78,12) {弱极限最优\\$J$ 弱序列\\下半连续};
\draw[ar] (s0) -- (a);
\draw[ar] (a) -- (b); \draw[ar] (b) -- (c); \draw[ar] (c) -- (d);
\node[fin,minimum width=48mm] (e) at (13,-8) {有界序列有收敛子列};
\node[fin] (f) at (52,-8) {极限仍可行\\$\bm{C}$ 闭};
\node[fin] (g) at (78,-8) {极限不跳高\\$f$ 下半连续};
\draw[VL,line width=.45pt,densely dashed] (a.south) -- (e.north -| a.south);
\draw[VL,line width=.45pt,densely dashed] (b.south) -- (e.north -| b.south);
\draw[VL,line width=.45pt,densely dashed] (c) -- (f);
\draw[VL,line width=.45pt,densely dashed] (d) -- (g);
\end{tikzpicture}
